\subsection{División de Líneas}

\subsubsection{Regla: cómo continuar una instrucción en varias líneas}

\textbf{Regla:} El estándar no fija una longitud máxima de línea. Cuando una instrucción se divide en varias líneas para que se lea mejor, se siguen estas convenciones:

\begin{itemize}
  \item un operador binario (\texttt{\&\&}, \texttt{||}, \texttt{??}, \texttt{+}) se coloca al inicio de la línea de continuación, no al final de la anterior;
  \item en una cadena de llamadas, cada \texttt{.Metodo()} va en su propia línea;
  \item si los argumentos de una llamada se dividen, la línea termina en el paréntesis de apertura y cada argumento va en su propia línea;
  \item la llamada \texttt{: base(...)} o \texttt{: this(...)} va en su propia línea, debajo de la firma del constructor;
  \item cada línea de continuación lleva un nivel más de indentación que la línea donde empieza la instrucción.
\end{itemize}

La opción \texttt{dotnet\_style\_operator\_placement\_when\_wrapping = beginning\_of\_line} aplica la primera convención en el editor (Microsoft, 2025e).

\textbf{Código correcto:}
\begin{lstlisting}[style=csharpstyle]
bool canMove = player.IsActive
    && board.IsInside(target)
    && !board.IsOccupied(target);

builder.Property(piece => piece.Kind)
    .HasMaxLength(KindLength)
    .IsUnicode(false);

throw Reject(
    new MoveLog(MoveResult.Blocked, attempt, player.PlayerId),
    MoveRejection.Blocked);
\end{lstlisting}

\textbf{Código incorrecto:}
\begin{lstlisting}[style=csharpstyle]
bool canMove = player.IsActive &&
    board.IsInside(target) &&
    !board.IsOccupied(target);

builder.Property(piece => piece.Kind).HasMaxLength(KindLength)
    .IsUnicode(false);

throw Reject(new MoveLog(MoveResult.Blocked, attempt,
    player.PlayerId), MoveRejection.Blocked);
\end{lstlisting}
